\newpage
\section{Tabellen, Symbole, Abkürzungen und Formeln}

\subsection{Tabelle}

Text\todo{Die folgenden Beispiele zeigen v.a. die Änderungen in der Formatierung, sodass die Arbeit den Vorgaben entspricht.} \par

\begin{table}[H]
\caption{Einfache Tabelle}
\begin{center}
\begin{tabularx}{0.7\textwidth}[ht]{|l|X|r|}
    \hline
Überschrift 1 & Überschrift 2 & Überschrift 3 \\ \hline
Inhalt 1 & Inhalt 2 & Inhalt 3 \\ \hline
1 & 2 & 3 \\ \hline
\end{tabularx}
\end{center}
\label{tbl:Beispieltabelle}
\end{table}
\vspace{-3em}
\cite[Quelle:][S. 300]{Beckert.2012}

\subsection{Symbole}

Anschleinend können Symbole folgenderweise verwendet werden: \\
Zuerst müssen sie in \texttt{skripte/symbolDef.tex} definiert werden. Ist das geschehen, können sie im Text referenziert und verwendet werden. \INF

\begin{lstlisting}
Zuerst müssen sie in \texttt{skripte/symbolDef.tex} definiert werden. Ist das geschehen, können sie im Text referenziert und verwendet werden. \INF
\end{lstlisting}

\subsection{Abkürzungen}

\ac{MS} Word verwendet den Ansatz \ac{WYSIWYG}. Wird eine Abkürzung mehrfach verwendet, sieht man, die Vorteile von einem Markup Format, das nicht nach \ac{WYSIWYG} arbeitet.

\subsection{Formeln}

Lorem ipsum dolor sit amet, consetetur sadipscing elitr, sed diam nonumy eirmod tempor invidunt ut labore et dolore magna aliquyam erat, sed diam voluptua. At vero eos et accusam et justo duo dolores et ea rebum. Stet clita kasd gubergren, no sea takimata sanctus est Lorem ipsum dolor sit amet. Lorem ipsum dolor sit amet, consetetur sadipscing elitr, sed diam nonumy eirmod tempor invidunt ut labore et dolore magna aliquyam erat, sed diam voluptua. At vero eos et accusam et justo duo dolores et ea rebum. Stet clita kasd gubergren, no sea takimata sanctus est Lorem ipsum dolor sit amet. Lorem ipsum dolor sit amet, consetetur sadipscing elitr, sed diam nonumy eirmod tempor invidunt ut labore et dolore magna aliquyam erat, sed diam voluptua. At vero eos et accusam et justo duo dolores et ea rebum. Stet clita kasd gubergren, no sea takimata sanctus est Lorem ipsum dolor sit amet.

\begin{xequation-}
	\caption{Momentanreserve bei einem kleinen Störfall}\label{Momentanreserve}
\begin{adjustwidth}{1cm}{0cm}
    \vspace{-1cm}
\begin{flalign}
	SIR_k &= \frac{RoCoF \cdot 2 \cdot W_k}{f_n} && \hfill(1) \nonumber \\
	\Leftrightarrow SIR_k &= \frac{0{,}125\,\frac{\mathrm{Hz}}{\mathrm{s}} \cdot 2 \cdot 4000\,\mathrm{MWs}}{50\,\mathrm{Hz}} && \hfill(1\text{a}) \nonumber \\
	\Leftrightarrow SIR_k &= 18\,\mathrm{MW}. && \hfill(1\text{b}) \nonumber
\end{flalign}
\end{adjustwidth}
\end{xequation-}
\cite[Quelle:][S. 420]{Wolff2017}

Lorem ipsum dolor sit amet, consetetur sadipscing elitr, sed diam nonumy eirmod tempor invidunt ut labore et dolore magna aliquyam erat, sed diam voluptua. At vero eos et accusam et justo duo dolores et ea rebum. Stet clita kasd gubergren, no sea takimata sanctus est Lorem ipsum dolor sit amet.

\subsection{Sammelwerk}

Wird ein Artikel aus einem Sammelwerk zitiert, funktioniert die Fußnote und das Literaturverzeichnis hoffentlich.~\footcite[\vglf][o. S.]{Liljegren2012}


\subsection{Glossar}

Brauchen wir in unserer Arbeit ein \Gls{glossar}? Falls ja, kann das mit dem Paket \glqq\gls{glossaries}\grqq\ umgesetzt werden.


